\section{理论深入：从 DDIM 到 DPM-Solver 的统一视角}

\subsection{DDIM 采样公式的重新推导}

讲者在课上带学生将 DPM-Solver-1 的结果逐项展开，验证其等价于 DDIM：

从 DPM-Solver-1 的公式出发：
$$\mathbf{x}_t = \frac{\alpha_t}{\alpha_s} \mathbf{x}_s - \sigma_t (e^{h} - 1) \boldsymbol{\epsilon}_\theta(\mathbf{x}_s, s)$$

其中 $h = \lambda_t - \lambda_s = \log\frac{\alpha_t}{\sigma_t} - \log\frac{\alpha_s}{\sigma_s} = \log\frac{\alpha_t \sigma_s}{\sigma_t \alpha_s}$。

因此 $e^h = \frac{\alpha_t \sigma_s}{\sigma_t \alpha_s}$，$e^h - 1 = \frac{\alpha_t \sigma_s - \sigma_t \alpha_s}{\sigma_t \alpha_s}$。

代入：
$$\mathbf{x}_t = \frac{\alpha_t}{\alpha_s} \mathbf{x}_s - \sigma_t \cdot \frac{\alpha_t \sigma_s - \sigma_t \alpha_s}{\sigma_t \alpha_s} \cdot \boldsymbol{\epsilon}_\theta(\mathbf{x}_s, s)$$
$$= \frac{\alpha_t}{\alpha_s} \mathbf{x}_s - \frac{\alpha_t \sigma_s - \sigma_t \alpha_s}{\alpha_s} \cdot \boldsymbol{\epsilon}_\theta(\mathbf{x}_s, s)$$

\begin{importantbox}{统一视角}
这正是 DDIM 在连续时间参数化下的采样公式。通过这一推导，我们看到了扩散模型不同理解方式之间的深层联系：
\begin{enumerate}
    \item DDPM $\to$ DDIM（Song et al., 2020）：从随机采样到确定性采样
    \item SDE $\to$ PF-ODE（Song et al., 2021）：从随机过程到确定性流
    \item PF-ODE $\to$ 指数积分（Lu et al., 2022）：从通用数值方法到结构化求解
\end{enumerate}
DDIM 恰好位于最后一条线的一阶近似位置——它不是"发明"出来的，而是扩散模型 ODE 结构的自然产物。
\end{importantbox}

\subsection{$\lambda$ 空间中的 Taylor 展开}

DPM-Solver 各阶方法的统一框架可以用 $\boldsymbol{\epsilon}_\theta$ 在 $\lambda$ 空间的 Taylor 展开来理解：
$$\boldsymbol{\epsilon}_\theta(\lambda) = \boldsymbol{\epsilon}_\theta(\lambda_s) + (\lambda - \lambda_s) \boldsymbol{\epsilon}_\theta'(\lambda_s) + \frac{(\lambda - \lambda_s)^2}{2} \boldsymbol{\epsilon}_\theta''(\lambda_s) + \cdots$$

\begin{itemize}
    \item \textbf{零阶截断}（$\boldsymbol{\epsilon}_\theta \approx$ 常数）$\to$ DPM-Solver-1 = DDIM
    \item \textbf{一阶截断}（$\boldsymbol{\epsilon}_\theta \approx$ 线性）$\to$ DPM-Solver-2
    \item \textbf{二阶截断}（$\boldsymbol{\epsilon}_\theta \approx$ 二次）$\to$ DPM-Solver-3
\end{itemize}

高阶导数 $\boldsymbol{\epsilon}_\theta'(\lambda_s)$ 无法解析获得（因为 $\boldsymbol{\epsilon}_\theta$ 是神经网络），但可以用有限差分来近似——这就是为什么需要在中间点额外评估网络。

\subsection{本章小结}

从 Taylor 展开的视角，DPM-Solver 的各阶方法形成了一个自然的精度递增族。DDIM 作为一阶特例的身份被揭示，高阶方法则是在 $\lambda$ 空间中对噪声预测做更精细的多项式拟合。


